\documentclass{article}
%Vsebina preambula.tex:
\usepackage[slovene]{babel}
\usepackage{lipsum} 
\usepackage{mathtools, amssymb, amsthm}
\usepackage{graphicx, caption, subcaption} 
\usepackage{array, booktabs}
\usepackage{xcolor}
\usepackage{hyperref}
\usepackage{ragged2e}

\NewDocumentCommand{\DemoColor}{m}{
    \raisebox{0.15cm}{#1} &
    \fcolorbox{black}{#1}{\phantom{\rule{0.5cm}{0.5cm}}}
  }

\NewDocumentCommand{\pozdrav}{O{svet}}
{Pozdravljeni, #1!}


\NewDocumentCommand{\kpi}{}
{\mathrm{\pi}} %Ukaz \kpi, kot 'konstantna pi'
\NewDocumentCommand{\ke}{}
{\mathrm{e}} %Ukaz \ke za Eulerjevo število


\NewDocumentCommand{\zap}{O{n}m}
{#2_{1}, \dots, #2_{#1}}

\NewDocumentCommand{\vektor}{om}
{
  \IfNoValueTF{#1}
  {(#2_{1}, #2_{2}, #2_{3})} %Če ni podan prvi argument, vrni vektor z 3 elementi
  {(#2_{1}, \dots, #2_{#1})} %Če je podan prvi argument, vrni vektor z #1 elementi
} 


\NewDocumentCommand{\ukaz}{o}
{
  \IfValueTF{#1}
  {Neobvezni argument je: \textbf{#1}} %Če je podan prvi argument, vrni njegov vrednost
  {Ni neobveznega argumenta.} %Če ni podanega argumenta
}

\NewDocumentCommand{\zapLim}{D<>{1}O{n}m}
{#3_{#1}, \dots, #3_{#2}}


\NewDocumentCommand{\jeRekli}{sm}
{\IfBooleanTF{#1}
{Ana je rekla: ``#2''}
{Boštjan je rekel: ``#2''}
}

\NewDocumentEnvironment{kralj}{}
{\begin{FlushLeft} \textbf{Poslušajte!} Njegova veličanstvo, kralj Artur bo podal izjavo: \end{FlushLeft}\begin{Center} \scshape  }
{\normalfont \end{Center}\begin{FlushRight} Konec izjave. Hvala za vašo pozornost. \end{FlushRight}}

\NewDocumentEnvironment{kraljIme}{m}
{\begin{FlushLeft} \textbf{Poslušajte!} Njegova veličanstvo, kralj \emph{#1} bo podal izjavo: \end{FlushLeft}\begin{Center} \scshape  }
{\normalfont \end{Center}\begin{FlushRight} Konec izjave. Hvala za vašo pozornost. \end{FlushRight}}

\NewDocumentEnvironment{kralji}{sm}
{ % ob \begin{kralji}
  \IfBooleanTF{#1}{ %z zvezdico
    \begin{FlushLeft} \textbf{Poslušajte!} Njegova veličanstvo, kralj \emph{#2} bo podal izjavo: \end{FlushLeft}\begin{Center} \scshape  
  }
  { %brez zvezdice
    \begin{FlushLeft} \textbf{Poslušajte!} Njena veličanstvo, kraljica \emph{#2} bo podala izjavo: \end{FlushLeft}\begin{Center} \scshape
  }
}
{ % ob \end{kralji}
  \IfBooleanTF{#1}
  { %z zvezdico
    \normalfont \end{Center}\begin{FlushRight} Konec izjave.  Kralj je hvalžen za vašo pozornost.\end{FlushRight}
  }
  { %brez zvezdice
    \normalfont \end{Center}\begin{FlushRight} Konec izjave. Kraljica je hvaležna za vašo pozornost. \end{FlushRight}
  }
}


\begin{document}
\color{red} Svet je zdaj rdeč, nekateri deli so \textcolor{yellow}{rumeni}. 
\[\mathcolor{blue}{\sum_{k=0}}^{10} i \]
Vse še vedno rdeča
\color{blue}
zdaj modra.


\begin{table}
  % \ExplSyntaxOn
  
  % \ExplSyntaxOff
  % \caption{Osnovne barve, ki so vnaprej določene v paketu \illa{xcolor}.}\label{tbl:basecolors}
  \begin{tabular}{@{}lc*2{@{\qquad}lc}@{}}
    \toprule
    Ime                 & Demo                  &
    Ime                 & Demo                  &
    Ime                 & Demo                                       \\
    \midrule
    \DemoColor{black}    & \DemoColor{lightgray} & \DemoColor{purple} \\
    \DemoColor{blue}     & \DemoColor{lime}      & \DemoColor{red}    \\
    \DemoColor{brown}    & \DemoColor{magenta}   & \DemoColor{teal}   \\
    \DemoColor{cyan}     & \DemoColor{olive}     & \DemoColor{violet} \\
    \DemoColor{darkgray} & \DemoColor{orange}    & \DemoColor{white}  \\
    \DemoColor{gray}     & \DemoColor{pink}      & \DemoColor{yellow} \\
    \DemoColor{green}    &                       &                    \\
    \bottomrule
  \end{tabular}
\end{table}

\vspace{2cm}

\textcolor{white!80!blue}{S}
\textcolor{white!60!blue}{l}
\textcolor{white!40!blue}{o}
\textcolor{white!20!blue}{v}
\textcolor{blue}{e}
\textcolor{blue!75!red}{n}
\textcolor{blue!50!red}{i}
\textcolor{blue!25!red}{j}
\textcolor{red}{a}


\vspace{2cm}
\noindent
\textcolor[Gray]{0}{Nula} \\
\textcolor[Gray]{3}{Tri} \\
\textcolor[Gray]{7}{Sedem} \\
\textcolor[Gray]{11}{Enaist} \\
\textcolor[Gray]{15}{Petnaist}

\vspace{2cm}


\definecolor{ULrdeca}{RGB}{226,25,30}
% V telesu dokumenta:
\noindent
\textcolor{ULrdeca}
{To je rdeča barva Univerze v Ljubljani.}

\color{black} % vrnemo barvo nazaj na črno

\newpage
% \pagecolor{yellow}
% \color{-yellow}
\noindent
Barva strani je rumena, \\
besedilo pa nasprotno rumeno.


\vspace{2cm}

\begin{minipage}{0.4\textwidth}
Zanimivo je, kako lahko barve \colorbox{yellow}{\textcolor{red}{izboljšajo}} dokument. Vendar pa lahko prekomerna uporaba barv dokument tudi \fcolorbox{orange}{gray}{pokvari}.
\end{minipage}

% \vspace{2cm}
% \begin{minipage}{0.4\textwidth}
%   Vsakič da uporabljamo ukaz \texttt{\textbackslash pozdrav},
%   dobimo isti rezultat: \pozdrav.
% \end{minipage}


\[
\ke^{i \kpi} + 1 = 0 \text{ vs. }
e^{i \pi} + 1 = 0
\]

\vspace{2cm}
\begin{minipage}{0.4\textwidth}
Dobrodošli na
  \textcolor[RGB]{226,25,30}{Univerzi v Ljubljani}!
  Ta \textcolor[RGB]{226,25,30}{barva} mi je zelo všeč.
\end{minipage}


\vspace{2cm}
\begin{minipage}{0.4\textwidth}
Ukaz \texttt{\textbackslash zap\{x\}\{n\}} ustvari zaporedje $x_{1}, \dots, x_{n}$,
torej \texttt{\textbackslash zap\{a\}\{10\}} daje $ \zap{a}{10} $.
\end{minipage}



\vspace{2cm}
\begin{minipage}{0.4\textwidth}
  Ukaz 
\texttt{\textbackslash pozdrav[<arg>]}
pozdravi \texttt{<arg>}, če ga ne damo, pozdravi \texttt{svet}:
\pozdrav[vsi], \pozdrav.
\end{minipage}

\vspace{2cm}
\begin{minipage}{0.4\textwidth}
Nekatere zaporedje imajo veliko elementov:
$\zap[1000]{a}$, 
za večino pa sploh ne vemo, 
koliko jih je:
$\zap{b}$.
\end{minipage}


\vspace{2cm}

\begin{minipage}{0.4\textwidth}
Na splošno ima vektor 
veliko vnosov: 
$\vektor[n]{x}$,
vendar večino časa delamo z 
vektorji s 3 vnosi:
  $\vektor{x}$.
\end{minipage}

\vspace{2cm}

\begin{minipage}{0.4\textwidth}
  \ukaz 
  \ukaz[Podan argument].
\end{minipage}

\vspace{2cm}

\begin{minipage}{0.4\textwidth}
  Večina zaporedj gre od $1$ do $n$:
  $\zapLim{a}$,
  vendar včasih želimo drugačen začetek:
  $\zapLim<0>{a}$,
  ali drugačen konec:
  $\zapLim[500]{a}$.
  ali pa kar oboje:
  $\zapLim<5>[500]{a}$.
\end{minipage}

\vspace{2cm}

\begin{minipage}
{0.4\textwidth}
  \jeRekli{Danes je lep dan.}\\
  \jeRekli*{Upam, da bo jutri tudi tako lep dan.}
\end{minipage}


\vspace{2cm}
\begin{kraljIme}{Artur}
  Kralj vam želi lepe praznike
\end{kraljIme}


\vspace{2cm}

\begin{kralji}*{Artur}
  Kraljeve besede.
\end{kralji}

\begin{kralji}{Guinevere}
  Kraljičine besede.
\end{kralji}




\end{document}